% Corrigé intégré automatiquement pour la banque Exercices.
% Source : A_Integrer/DDS_03/048_PFS/048_PFS.tex
\begin{corrigebox}[Corrigé]
\CorrigeQuestion{1}
La résolution consiste à identifier les données et l'inconnue, choisir la loi du cours adaptée, écrire l'équation symbolique avant toute application numérique, puis vérifier l'homogénéité, le signe et l'ordre de grandeur du résultat. Les hypothèses utilisées doivent être explicitement contrôlées à la fin.
\CorrigeQuestion{2}
On isole le solide ou l'ensemble indiqué, on dresse le bilan complet des actions mécaniques extérieures et on choisit un point de réduction qui élimine le maximum d'inconnues. Le PFS donne $\sum\vec F=\vec0$ et $\sum\vec M_A=\vec0$ ; les projections utiles fournissent les inconnues, puis leur signe permet de vérifier le sens supposé.
\CorrigeQuestion{3}
On isole le solide ou l'ensemble indiqué, on dresse le bilan complet des actions mécaniques extérieures et on choisit un point de réduction qui élimine le maximum d'inconnues. Le PFS donne $\sum\vec F=\vec0$ et $\sum\vec M_A=\vec0$ ; les projections utiles fournissent les inconnues, puis leur signe permet de vérifier le sens supposé.
\end{corrigebox}
